
\section{Kết luận}

Báo cáo đã phân tích và tái lập hai công trình
CateEA và DSoM, lần lượt giải quyết bài toán
căn chỉnh thực thể đa phương thức (MMEA)
và suy luận trên đồ thị tri thức đa phương thức
(MKGR).

Hai công trình có cách tiếp cận khác nhau.
CateEA khai thác thông tin danh mục ngầm
thông qua phân cụm để bổ sung tín hiệu
giám sát cho quá trình căn chỉnh thực thể.
Trong khi đó, DSoM sử dụng cơ chế lựa chọn
teacher và Knowledge Distillation để truyền
tri thức đa phương thức sang student
chỉ sử dụng cấu trúc.

Kết quả tái lập phần lớn phản ánh các xu hướng
được công bố, đồng thời cho thấy một số
sai khác giữa phương pháp và mã nguồn.
Đối với CateEA, nhóm ghi nhận những vấn đề
liên quan đến cách gán nhãn giả và điều chỉnh
độ tương đồng. Đối với DSoM, nhóm phát hiện
policy RC không được cập nhật trong cấu hình
mã nguồn công khai.

Qua hai công trình, nhóm nhận thấy việc
đề xuất một cơ chế học mới không đồng nghĩa
với việc đóng góp của từng thành phần
đã được kiểm chứng đầy đủ. Các thực nghiệm
ablation, việc đối chiếu mã nguồn và đánh giá
trên nhiều random seed là cần thiết để
xác định mức độ hiệu quả và tính ổn định
của phương pháp.

Nhìn chung, quá trình phân tích và tái lập
không chỉ giúp nhóm hiểu rõ hơn hai phương pháp mà còn làm rõ vai trò của việc kiểm chứng
thực nghiệm khi đánh giá một công trình
nghiên cứu về học máy trên dữ liệu đồ thị.
